Este relatório é a primeira versão do relatório final do sistema de monitoramento não
invasivo de idosos em Instituições de Longa Permanência (ILPIs). Em cada quarto, um
\textbf{hub} com ESP32-S3 lê as portas (reed switch MC-38), a presença (PIR HC-SR501) e a
temperatura (LM35DZ condicionado por um LM358), roda as regras de alerta e publica tudo
por MQTT. O idoso usa um \textbf{vestível}, também com ESP32-S3, com um acelerômetro e
giroscópio MPU-6050 a 200\,Hz, que detecta quedas com um Random Forest embarcado e avisa
o hub por ESP-NOW cifrado. Uma \textbf{central} em Docker guarda o histórico (InfluxDB),
mostra os quartos (Grafana) e avisa o cuidador (Telegram).

As duas placas são a mesma \textbf{PCB de duas camadas}, desenhada no KiCad e fabricada
na JLCPCB, em que a DevKitC-1 é encaixada; só o firmware muda. O documento cobre os
sensores e a caracterização deles, o condicionamento do LM35 (ganho de 5,5 em
alimentação simples e filtro RC de \rcFc\,Hz, com a curva
$T = (V - \lmBetaCadeia)/\lmAlfaCadeia$ ajustada em \lmPontos{} pontos), a placa, a
interface CAN (TWAI a 250\,kbit/s, identificadores \code{0x0GS}), a interface sem fio
(ESP-NOW do vestível ao hub, MQTT do hub à central), o firmware em FreeRTOS, a detecção de
queda e a central. Na base SisFall, o Random Forest do vestível dá \rfAlarmes{} alarme
falso por hora de atividade e detecta \rfQuedas\,\% das quedas; o modelo ainda é
provisório, porque foi treinado com o sensor na cintura.

O firmware das duas placas compila sem avisos e a lógica pura passa em mais de mil
verificações no computador, mas o sistema completo ainda não rodou em bancada: a
Seção~\ref{sec:verificacao} separa o que foi verificado, e como, do plano de testes com
as placas montadas.

{\raggedright
\textbf{Palavras-chave:} ILPI; detecção de queda; ESP32-S3; rede CAN; ESP-NOW; MQTT;
condicionamento de sinais; PCB; FreeRTOS.\par}
